\section{Learning to Desire What Heaven Gives}

Holiness prepares us for the actual happiness of heaven. The blessed love God without reserve; they love one another in him; they rejoice in truth. A life that deliberately shelters hatred, dishonesty, or impurity is therefore moving against the good it professes to seek. Hatred prevents us from sharing another's joy, and dishonesty makes truth unwelcome. Our capacity to receive God's gift needs healing. The Lord who forgives also makes holy, and the cleansing he completes in those who die in his friendship belongs to the generosity with which he brings them into his presence.

In the sermon ``Holiness Necessary for Future Blessedness,'' preached before his reception into the Catholic Church, John Henry Newman considers this relation between the heart's character and its happiness (\emph{Parochial and Plain Sermons} I.1). A person does not become able to delight in God's holiness merely by entering a different place. Heaven is life in God's presence; the heart must be made capable of loving that presence. Newman's imagined case of an unholy person admitted unchanged is a spiritual warning about the need for conversion, not an account of someone actually unhappy within the beatific vision. The blessed are holy because grace has made them so.

Newman compares this preparation to learning to take part in worship. In prayer we give attention to God, confess our need, praise him, and seek his will. These acts have their proper end in a heart turned toward him. Prayer can expose how much of our attention we are willing to give to God. We may find it easy to spend ourselves on an ambition and difficult to give God a few attentive minutes. The discovery should bring us to him honestly. It need not be converted into a verdict that every distracted or dry prayer proves a lack of love. Fidelity in dryness can itself be an act of love.

Prayer also exposes the conditions we set upon our obedience. We can practice an outward duty while preserving a hidden condition: we will serve so long as we are appreciated, forgive so long as the other person remains beneath us, or pray so long as prayer offers consolation. God patiently brings these conditions into the light. Their surrender enlarges our capacity for charity. A person who accepts correction and returns to a duty without the reassurance he wanted may be growing toward heaven in a way that remains invisible to him. He may recognize the growth only later, in a changed response to the same difficulty.

Newman's sermon ends by directing this demanding work toward the help of the Holy Spirit. God is at work within us, and his help makes the change possible. Without that assurance, a more truthful knowledge of our attachments can become discouragement. With it, recognizing an attachment gives us something definite to bring to God for healing. We can ask for a desire we do not yet fully possess, and we can act upon the grace given today without first feeling the completion that belongs to glory.

St.~Francis de Sales makes the same movement practical in his meditation on paradise (\emph{Introduction to the Devout Life} I.16). He begins with beauty, draws the reader toward the company of the saints, and brings the consideration to God's presence. The affections of admiration and desire then become a resolution to relinquish particular hindrances and undertake particular helps. The resolution answers the desire awakened by the meditation. A particular attachment or neglected help can now be considered in relation to the good Francis has asked the reader to contemplate.

A renewed desire for heaven may therefore take a very ordinary form. It may mean putting aside a source of vanity, being faithful to a neglected prayer, or ending the deliberate rehearsal of resentment. A resolution can be kept when it names something we can actually do and a hindrance we can recognize. We do not know how much purification remains. We do know that every consent to grace belongs to the journey he has opened. We welcome that purification because we hope to see him face to face.

